\documentclass[11pt]{article}
\usepackage[margin=1in]{geometry}
\usepackage{amsmath,amssymb,amsthm}
\usepackage{algorithm}
\usepackage{algpseudocode}
\usepackage{booktabs}
\usepackage{hyperref}

\newtheorem{definition}{Definition}
\newtheorem{theorem}{Theorem}
\newtheorem{lemma}[theorem]{Lemma}
\newtheorem{property}{Property}

\title{\textbf{Mathematical Specification, Dual Number Algebra, and Directional Tangent Propagation of Forward-Mode Differentiation (ALGO-NN-25)}}
\author{Team Scaibu \\ \small Email: \texttt{jaydeep.9664920749@gmail.com} \\ \small Architecture and Core Engine Team}
\date{\today}

\begin{document}
\maketitle

\begin{abstract}
Forward-mode automatic differentiation propagates primal values and directional tangent perturbations simultaneously using the commutative ring of dual numbers $\mathbb{R}[\epsilon]/(\epsilon^2)$. For functions $f: \mathbb{R}^n \to \mathbb{R}^m$ where $m \gg n$, forward AD evaluates exact Jacobian-Vector Products (JVPs) in a single forward pass without requiring activation storage tapes. We present the formal algebraic axioms, prove the JVP Equivalence Theorem, provide complete pseudocode matching the reference implementation, and derive complexity bounds.
\end{abstract}

\section{Notation and Mathematical Preliminaries}
\label{sec:notation}

Let $\mathbb{D} = \{a + b\epsilon \mid a, b \in \mathbb{R}, \epsilon^2 = 0, \epsilon \neq 0\}$ denote the ring of dual numbers. For any differentiable mapping $f: \mathbb{R}^n \to \mathbb{R}^m$, let $\mathbf{x} \in \mathbb{R}^n$ represent the primal point and $\mathbf{v} \in \mathbb{R}^n$ the direction tangent vector.

\subsection{Dual Number Arithmetic Mechanics}
\paragraph{Intuitive Meaning:} By defining $\epsilon$ as an infinitesimal number whose square is zero, evaluating a function at $x + \epsilon v$ yields the standard function output in the real part and the exact directional derivative in the dual part.

\paragraph{Formal Mathematical Definition:}
Applying Taylor series expansion about primal $x$:
\begin{equation}
\label{eq:dual_taylor}
f(x + \epsilon v) = f(x) + f'(x) \epsilon v + \frac{1}{2} f''(x) (\epsilon v)^2 + \dots = f(x) + \epsilon (f'(x) v).
\end{equation}

\paragraph{Worked Numerical Example:}
Let $f(x) = x^2 + \sin(x)$ at $x = 2.0$ with tangent $v = 1.0$:
Primal: $2.0^2 + \sin(2.0) = 4.0 + 0.909297 = 4.909297$.
Tangent: $(2(2.0) + \cos(2.0)) \times 1.0 = (4.0 - 0.416147) \times 1.0 = 3.583853$.

\subsection{Summary of Symbols}
\begin{itemize}
    \item $N \in \mathbb{N}$: Number of primal/tangent input variables.
    \item $\mathbf{x} \in \mathbb{R}^N$: Primal input vector.
    \item $\mathbf{v} \in \mathbb{R}^N$: Tangent perturbation vector.
    \item $\mathbf{J} \in \mathbb{R}^{M \times N}$: Underlying Jacobian matrix $\left[\frac{\partial f_i}{\partial x_j}\right]$.
    \item $\mathbf{J}\mathbf{v} \in \mathbb{R}^M$: Evaluated Jacobian-Vector Product.
\end{itemize}

\section{Problem Definition and Core Concepts}
\label{sec:problem_definition}

\subsection{Elementary Dual Operations}
\begin{itemize}
    \item \textbf{Square ($y = x^2$):} $y_{\mathrm{primal}} = x_p^2, \quad y_{\mathrm{tangent}} = 2 x_p x_t$.
    \item \textbf{Sine ($y = \sin(x)$):} $y_{\mathrm{primal}} = \sin(x_p), \quad y_{\mathrm{tangent}} = \cos(x_p) x_t$.
    \item \textbf{Exponential ($y = \exp(x)$):} $y_{\mathrm{primal}} = \exp(x_p), \quad y_{\mathrm{tangent}} = y_p x_t$.
    \item \textbf{ReLU ($y = \max(0, x)$):} $y_{\mathrm{primal}} = \max(0, x_p), \quad y_{\mathrm{tangent}} = x_t \cdot \mathbb{I}_{\{x_p > 0\}}$.
    \item \textbf{Linear Combination ($y = \sum w_k x_k$):} $y_p = \sum w_k x_{p, k}, \quad y_t = \sum w_k x_{t, k}$.
\end{itemize}

\section{Algorithmic Specification}
\label{sec:algorithm}

Algorithm~\ref{alg:forward_ad} specifies the primal-tangent forward pass matching \texttt{impl.py}.

\begin{algorithm}[H]
\caption{Forward-Mode Differentiation and Jacobian-Vector Product Evaluation}
\label{alg:forward_ad}
\begin{algorithmic}[1]
\Require Primal inputs $\mathbf{x} \in \mathbb{R}^N$, Tangent inputs $\mathbf{v} \in \mathbb{R}^N$, operation sequence $\mathcal{O} = [O_1, \dots, O_K]$.
\Ensure Evaluated primals $\mathbf{y}_{\mathrm{primal}}$, evaluated tangents $\mathbf{y}_{\mathrm{tangent}} = \mathbf{J}\mathbf{v}$.
\State Initialize primals $\mathbf{p} \gets \text{list}(\mathbf{x}), \quad \mathbf{t} \gets \text{list}(\mathbf{v})$
\For{each operation $O_k \in \mathcal{O}$}
    \State $\mathrm{op} \gets O_k.\mathrm{op}, \quad \mathrm{parents} \gets O_k.\mathrm{parents}, \quad \mathbf{w} \gets O_k.\mathrm{weights}$
    \State $p_0 \gets p_{\mathrm{parents}[0]}, \quad t_0 \gets t_{\mathrm{parents}[0]}$
    \If{$\mathrm{op} = \text{square}$}
        \State $p_{\mathrm{new}} \gets p_0^2, \quad t_{\mathrm{new}} \gets 2.0 \cdot p_0 \cdot t_0$
    \ElsIf{$\mathrm{op} = \text{sin}$}
        \State $p_{\mathrm{new}} \gets \sin(p_0), \quad t_{\mathrm{new}} \gets \cos(p_0) \cdot t_0$
    \ElsIf{$\mathrm{op} = \text{exp}$}
        \State $p_{\mathrm{new}} \gets \exp(p_0), \quad t_{\mathrm{new}} \gets p_{\mathrm{new}} \cdot t_0$
    \ElsIf{$\mathrm{op} = \text{relu}$}
        \State $p_{\mathrm{new}} \gets \max(0.0, p_0), \quad t_{\mathrm{new}} \gets t_0 \cdot \mathbb{I}_{\{p_0 > 0\}}$
    \ElsIf{$\mathrm{op} = \text{linear\_combination}$}
        \State $p_{\mathrm{new}} \gets \sum_j w_j p_{\mathrm{parents}[j]}, \quad t_{\mathrm{new}} \gets \sum_j w_j t_{\mathrm{parents}[j]}$
    \EndIf
    \State Append $p_{\mathrm{new}}$ to $\mathbf{p}$, append $t_{\mathrm{new}}$ to $\mathbf{t}$
\EndFor
\State \Return $\mathbf{p}, \quad \mathbf{t}$
\end{algorithmic}
\end{algorithm}

\section{Theoretical Guarantees and JVP Equivalence}
\label{sec:guarantees}

\begin{theorem}[Exact Jacobian-Vector Product Equivalence]
\label{thm:jvp_exactness}
\textbf{Assumptions:} Let $f: \mathbb{R}^N \to \mathbb{R}^M$ be composed of elementary differentiable operations, evaluated at primal $\mathbf{x}$ with tangent direction $\mathbf{v}$.
\textbf{Guarantee:} Forward-mode dual number propagation evaluates the exact directional derivative $\mathbf{J} \mathbf{v} = \lim_{\epsilon \to 0} \frac{f(\mathbf{x} + \epsilon \mathbf{v}) - f(\mathbf{x})}{\epsilon}$ in $\operatorname{Time}(\text{JVP}) \le 2.5 \operatorname{Time}(f)$ with $\mathcal{O}(1)$ working memory overhead beyond forward pass storage.
\textbf{Proof:}
Every elementary unary/binary dual operation $(a + \epsilon \dot{a}) \circ (b + \epsilon \dot{b})$ computes $(a \circ b) + \epsilon (\frac{\partial \circ}{\partial a} \dot{a} + \frac{\partial \circ}{\partial b} \dot{b})$ by definition of dual arithmetic. By mathematical induction over DAG depth, each node $k$ carries $v_k = f_k(\mathbf{x})$ and $\dot{v}_k = \sum_j \frac{\partial f_k}{\partial x_j} v_j = (\mathbf{J}_k \mathbf{v})$. Thus the final outputs contain the exact JVP without finite-difference truncation error.
\textbf{Limitations:} Evaluating full Jacobian matrices $\mathbf{J} \in \mathbb{R}^{M \times N}$ requires $N$ sequential forward passes, making reverse-mode AD preferred when $N \gg M$ (as in scalar loss optimization).
\end{theorem}

\section{Complexity Analysis}
\label{sec:complexity}

\begin{itemize}
    \item \textbf{Time Complexity:} $\mathcal{O}(M)$ operations for forward pass with constant factor $\le 2.5\times$ pure primal evaluation.
    \item \textbf{Space Complexity:} $\mathcal{O}(M)$ online storage, with zero activation graph retention needed after node computation.
\end{itemize}

\section{Worked Numerical Example}
\label{sec:worked_example}

Let inputs $x_0 = 3.0, v_0 = 1.0$:
Operation 1: $x_1 = \text{square}(x_0) \implies p_1 = 9.0, t_1 = 2(3.0)(1.0) = 6.0$.
Operation 2: $x_2 = \sin(x_0) \implies p_2 = \sin(3.0) = 0.141120, t_2 = \cos(3.0)(1.0) = -0.989992$.
Operation 3: $x_3 = x_1 + x_2 \implies p_3 = 9.141120, t_3 = 6.0 - 0.989992 = 5.010008$.

\section{Numerical Considerations and Edge Cases}
\label{sec:numerical}

\begin{itemize}
    \item \textbf{Memory Footprint:} Unlike reverse-mode AD, forward AD does not save intermediate layer states, allowing evaluation on arbitrarily deep networks under fixed memory bounds.
\end{itemize}

\begin{thebibliography}{9}
\bibitem{wengert1964} R.~E.~Wengert, ``A Simple Automatic Derivative Evaluation Program,'' \emph{Communications of the ACM}, vol.~7, no.~8, pp.~463--464, 1964.
\end{thebibliography}

\end{document}
